% !TeX program = xelatex
\documentclass[12pt,a4paper]{article}
\usepackage[margin=2cm]{geometry}
\usepackage{fontspec}
\usepackage{amsmath,amssymb}
\setmainfont{Times New Roman}
\setlength{\parindent}{0pt}
\setlength{\parskip}{6pt}
\begin{document}
\begin{center}\Large\bfseries Đáp án Toán vào lớp 10 Hà Nội năm 2012\\Phần Đại số\end{center}

\textbf{Câu 16.}

a) Thay $x=36$:
\[
A=\frac{\sqrt{36}+4}{\sqrt{36}+2}=\frac{6+4}{6+2}=\frac54.
\]

b) Đặt $t=\sqrt{x}$. Khi đó $t\ge0$, $t\ne4$, và $x=t^2$. Ta có
\begin{align*}
B&=\left(\frac{t}{t+4}+\frac4{t-4}\right):\frac{t^2+16}{t+2}\\
&=\frac{t(t-4)+4(t+4)}{(t+4)(t-4)}\cdot\frac{t+2}{t^2+16}\\
&=\frac{t^2+16}{(t+4)(t-4)}\cdot\frac{t+2}{t^2+16}\\
&=\frac{t+2}{t^2-16}=\frac{\sqrt{x}+2}{x-16}.
\end{align*}
Điều kiện ban đầu đảm bảo các mẫu khác $0$.

c) Từ hai kết quả trên:
\[
A-1=\frac{2}{\sqrt{x}+2},\qquad B(A-1)=\frac{2}{x-16}.
\]
Biểu thức này là số nguyên khi $x$ nguyên (ngoài điều kiện $x\ge0$, $x\ne16$) và
\[
\frac2{x-16}\in\mathbb Z
\iff x-16\in\{-2,-1,1,2\}.
\]
Do đó $x\in\{14,15,17,18\}$.

\textbf{Câu 17.} Gọi $x,y$ lần lượt là số giờ người thứ nhất và người thứ hai làm một mình xong việc. Khi đó $y-x=2$, $x>0$. Năng suất mỗi giờ của hai người lần lượt là $\dfrac1x$ và $\dfrac1y$. Vì cùng làm thì hoàn thành trong $\dfrac{12}{5}$ giờ:
\[
\frac{12}{5}\left(\frac1x+\frac1y\right)=1.
\]
Thay $y=x+2$:
\[
\frac{12}{5}\left(\frac1x+\frac1{x+2}\right)=1
\iff 5x^2-14x-24=0
\iff (x-4)(5x+6)=0.
\]
Vì $x>0$, suy ra $x=4$, nên $y=6$. Người thứ nhất cần $4$ giờ và người thứ hai cần $6$ giờ.

\textbf{Câu 18.}

a) Đặt $u=\dfrac1x$, $v=\dfrac1y$. Hệ trở thành
\[
\begin{cases}2u+v=2,\\6u-2v=1.\end{cases}
\]
Nhân phương trình đầu với $2$ rồi cộng phương trình sau được $10u=5$, vậy $u=\dfrac12$ và từ đó $v=1$. Suy ra $x=2$, $y=1$. Nghiệm là $(x;y)=(2;1)$.

b) Biệt thức của phương trình theo $x$ là
\[
\Delta=(4m-1)^2-4(3m^2-2m)=4m^2+1>0
\]
với mọi $m$, nên phương trình luôn có hai nghiệm phân biệt. Theo Viète,
\[
x_1+x_2=4m-1,\qquad x_1x_2=3m^2-2m.
\]
Điều kiện đề bài tương đương
\begin{align*}
7=x_1^2+x_2^2&=(x_1+x_2)^2-2x_1x_2\\
&=(4m-1)^2-2(3m^2-2m)\\
&=10m^2-4m+1,
\end{align*}
nên $10m^2-4m-6=0$, hay $(m-1)(5m+3)=0$. Vậy $m=1$ hoặc $m=-\dfrac35$.

\textbf{Câu 20.} Đặt $t=\dfrac{x}{y}$. Vì $x\ge2y$ và $y>0$ nên $t\ge2$. Khi đó
\[
M=\frac{x^2+y^2}{xy}=t+\frac1t.
\]
Ta có
\[
M-\frac52=\frac{2t^2-5t+2}{2t}
=\frac{(t-2)(2t-1)}{2t}\ge0,
\]
do $t\ge2$. Vậy $M_{\min}=\dfrac52$, đạt được khi $t=2$, tức $x=2y$.
\end{document}
